\documentclass[11pt]{article}

\usepackage{amsmath, amssymb, amsthm, amsfonts}
\usepackage[hidelinks]{hyperref}
\usepackage{geometry}
\geometry{a4paper, margin=1in}

\theoremstyle{definition}
\newtheorem{definition}{Definition}

\theoremstyle{plain}
\newtheorem{proposition}{Proposition}

\theoremstyle{remark}
\newtheorem{example}{Example}

\title{Explanation of the Invariant $r = \alpha x - \beta y$\\
\large and Arithmetic Progressions in the Proof of Conjecture~3}
\author{Dirk Kunert \thanks{Independent Researcher, Email: dirk.kunert@gmail.com}
\and
Claude Opus 4.6\thanks{AI research assistant (Anthropic).}}
\date{\today}

\begin{document}

\maketitle

% =====================================================================
\section{Arithmetic Progressions}
% =====================================================================

\begin{definition}
An \emph{arithmetic progression} is a sequence of numbers in which the
difference between consecutive terms is constant.  This constant is called
the \emph{common difference}.
\end{definition}

\noindent\textbf{General form.}
An arithmetic progression with initial value $a$ and common difference $d$ is
the sequence
\[
    a,\; a + d,\; a + 2d,\; a + 3d,\; \ldots
\]
or, as a set indexed by $k \in \mathbb{Z}$:
\[
    \{\, a + kd \;:\; k \in \mathbb{Z} \,\}.
\]

\begin{example}
The sequence $3, 7, 11, 15, 19, \ldots$ is an arithmetic progression with
initial value $a = 3$ and common difference $d = 4$.
\end{example}

\noindent\textbf{Role in the proof.}
In the proof of Conjecture~3, each set
\[
    P_r \;=\; \bigl\{\, c_r + k(\alpha^2 + \beta^2) \;:\; k \in \mathbb{Z} \,\bigr\}
\]
is an arithmetic progression with initial residue $c_r$ and common difference
$D = \alpha^2 + \beta^2$.  The projected values $\tilde{p} = \beta x + \alpha y$
for a fixed value of $r$ are therefore equally spaced with spacing
$\alpha^2 + \beta^2$.

There are $\alpha + \beta + 1$ such progressions (one for each value of $r$),
all sharing the same common difference but having different offsets $c_r$.
The interleaving of these progressions --- sorting all their elements together
--- produces the periodic difference sequence whose period length is
$\alpha + \beta + 1$.

% =====================================================================
\section{Why $r = \alpha x - \beta y$ Is Called an Invariant}
% =====================================================================

The quantity $r := \alpha x - \beta y$ is called an \emph{invariant} because
it remains constant as one moves along a given arithmetic progression of
lattice points.

\begin{proposition}
Within a fixed progression $P_r$, the lattice points are parametrised by
$k \in \mathbb{Z}$ as
\[
    x \;=\; x_0 + \beta k, \qquad y \;=\; y_0 + \alpha k,
\]
where $(x_0, y_0)$ is a base point with $\alpha x_0 - \beta y_0 = r$.
Then $\alpha x - \beta y = r$ for every $k$.
\end{proposition}

\begin{proof}
Substituting the parametrisation:
\[
    \alpha(x_0 + \beta k) - \beta(y_0 + \alpha k)
    \;=\; \alpha x_0 - \beta y_0
        + \underbrace{\alpha \beta k - \beta \alpha k}_{=\, 0}
    \;=\; r.
\]
The $k$-dependent terms cancel exactly because the step direction
$(\beta, \alpha)$ is the direction of the line
$f(x) = \frac{\alpha}{\beta}\,x$.
\end{proof}

\noindent\textbf{Geometric interpretation.}
The line $f(x) = \frac{\alpha}{\beta}\,x$ has direction vector
$(\beta, \alpha)$.  Moving along this direction changes $x$ by $\beta$ and
$y$ by $\alpha$, leaving $\alpha x - \beta y$ unchanged.  The quantity $r$
therefore labels which of the $\alpha + \beta + 1$ parallel ``tracks''
a lattice point belongs to.  Points on the same track share the same $r$
and lie on the same arithmetic progression of $\tilde{p}$-values.
Points on different tracks have different values of $r$.

\begin{example}
For $\alpha = 2$, $\beta = 1$, and $\omega = 1$, the invariant takes values
in $\mathcal{R} = \{-1, 0, 1, 2\}$, giving $|\mathcal{R}| = 4 = 2 + 1 + 1
= \Lambda_{2,1}$.
\begin{itemize}
    \item $r = -1$: e.g.\ $(x, y) = (1, 3)$, since $2 \cdot 1 - 1 \cdot 3 = -1$.
          Next point on same track: $(1 + 1, 3 + 2) = (2, 5)$, and indeed
          $2 \cdot 2 - 1 \cdot 5 = -1$.
    \item $r = 0$: e.g.\ $(x, y) = (1, 2)$, since $2 \cdot 1 - 1 \cdot 2 = 0$.
    \item $r = 1$: e.g.\ $(x, y) = (1, 1)$, since $2 \cdot 1 - 1 \cdot 1 = 1$.
    \item $r = 2$: e.g.\ $(x, y) = (1, 0)$, since $2 \cdot 1 - 1 \cdot 0 = 2$.
\end{itemize}
Each of these four tracks produces its own arithmetic progression of
$\tilde{p}$-values with common difference $D = 4 + 1 = 5$.
Their interleaving gives the difference sequence $(1, 1, 1, 2)$ with period
$4 = \Lambda_{2,1}$.
\end{example}

\end{document}
